% ~300 words.
\section{Discussion}

FLAIR combines three properties that, individually, have appeared
elsewhere but together have not been operationalized at the scale ICU
CDS development requires: federation across many heterogeneous sites
without patient-data movement, a clinically scoped data schema (CLIF)
that makes the analytic code portable, and an extension protocol that
lets the task catalog grow by community pull request rather than by
single-author fiat.

\paragraph{Saturation resistance through siloed data.} A side effect of
the model-to-data design is that the data underlying FLAIR's leaderboard
never enters public hands. Test-set contamination of pretraining
corpora --- now a recurring problem for open benchmarks
\citep{recht2019imagenet,chekroud2024illusory} --- is structurally
limited, and the incentive for a single developer to overfit the
leaderboard is weak because the data they would overfit to are not
accessible to them. We view this not as a feature of FLAIR per se but
as an emergent property of the access model.

\paragraph{Limitations.} Adoption is bounded by the institutions that
have implemented CLIF; non-CLIF sites cannot currently participate
without an ETL effort. The current task suite is restricted to binary
classification on structured EHR variables and excludes notes,
waveforms, and images. Aggregate-only outputs constrain some
fine-grained bias analyses; subgroups below the suppression threshold
are not visible. T5 currently uses serum-creatinine-only KDIGO, and
T4's antibiotic vocabulary is an open dependency.

\paragraph{Roadmap.} Near-term work extends the suite to opt-in
prospective silent-trial evaluation \citep{joshi2025silenttrials},
longitudinal calibration drift monitoring against site reference
periods, and integration of additional CLIF tables (diagnoses,
procedures) as the schema matures. The model-to-data primitive
underpinning FLAIR is data-modality-agnostic; once CLIF stabilizes
multimodal extensions, the same submission and aggregation contract
will apply to notes and waveforms.

\paragraph{Conclusion.} FLAIR makes privacy-preserving multi-site
external validation of EHR-based clinical decision support a routine,
reproducible operation. The library, task definitions, container, and
report generators are open source (Apache~2.0,
\texttt{github.com/clif-consortium/FLAIR}); researchers and clinicians
are invited to add new benchmark tasks through the single-file PR
template and two-site governance review described in
\S\ref{sec:benchmarks}, so the catalog grows with the field.
